\chapter*{Abstract}
\addcontentsline{toc}{chapter}{Abstract}

This work investigates the application of Abaqus built-in error-guided mesh refinement and adaptive remeshing capabilities to phase-field fracture simulations implemented with user elements. The point of departure is the staggered phase-field implementation of Moln\'ar and Gravouil and the mesh-refinement procedure proposed by Pandey and Kumar. The core research objectives comprise: (1) reproducing reference fracture simulations without mesh refinement; (2) implementing the Abaqus-native refinement procedure in Python while retaining the layered UEL/UMAT/facsimile representation required by the user-element formulation; (3) quantitatively reproducing the reference results with adaptive refinement; (4) integrating and verifying the IMFD ABAQUSER post-processing tool; (5) applying the adaptive workflow to fracture benchmarks and sensitivity studies; and (6) assessing nonmatching state transfer and restart mechanics across evolving discretizations.

First, a fixed-mesh Mode-I single-edge-crack benchmark was reproduced without mesh refinement (Job \texttt{1398090.mmaster02}, $15{,}192$ Quad4 elements, walltime $06\,\mathrm{h}\,31\,\mathrm{min}$, CPUT $06\,\mathrm{h}\,30\,\mathrm{min}$). The reconstructed response achieved a peak reaction force of $F_{\text{peak}} = 0.757778\,\mathrm{kN}$ at $u_{\text{peak}} = 0.005857\,\mathrm{mm}$ compared to the published target of $0.7580\,\mathrm{kN}$ at $0.005860\,\mathrm{mm}$, exhibiting relative differences of only $-0.029\%$ in peak load and $-0.051\%$ in peak displacement. This fixed-mesh calculation established the quantitative baseline for the refinement study (Task 3: \textbf{COMPLETE / SCIENTIFICALLY ACCEPTED}).

Second, an Abaqus-native Python-controlled two-job pre-refinement workflow was implemented following Pandey and Kumar. A coarse pre-analysis supplies the Abaqus recovered-stress discretization error indicator \texttt{MISESERI}; an Abaqus \texttt{RemeshingRule} and \texttt{adaptiveRemesh} operation generate the refined physical mesh; and the resulting connectivity is converted automatically into coincident phase-field UEL, mechanical UEL, and visualization UMAT layers. Qualification testing demonstrated genuine native remeshing and eliminated an implementation defect where calling \texttt{Part.generateMesh()} after \texttt{adaptiveRemesh()} had previously overwritten the refined mesh with coarse part seeds (Task 4: \textbf{COMPLETE / VERIFIED END-TO-END}).

Third, production adaptive Mode-I simulations were performed and reconciled. Reconstructing the nominal $1.0\%$ error target reported in the literature generated an over-refined mesh of $71{,}320$ physical elements ($69{,}443$ Quad4, $1{,}877$ Tri3) on a sharp crack seam, which completed technically (Job \texttt{1399632.mmaster02}, walltime $35\,\mathrm{h}\,08\,\mathrm{min}$, CPUT $34\,\mathrm{h}\,03\,\mathrm{min}$) but failed quantitative reproduction gates due to an elastic compliance shift and accelerated crack-tip damage localization ($F_{\text{peak}} = 0.478203\,\mathrm{kN}$, $-36.91\%$). A controlled forensic audit established that because the relative stress error $\eta = \text{MISESERI}/\text{MISESAVG}$ is scale-invariant across a singular crack tip, the nominal $1.0\%$ threshold marks $>80\%$ of the domain for refinement in our native implementation. While this mechanism explains our observed over-refinement, the exact cause of why the publication reported $13{,}941$ elements cannot be uniquely determined from published information and is classified as \textbf{\texttt{DOCUMENTED\_UNRESOLVABLE\_FROM\_PUBLISHED\_INFORMATION}}. An empirical $2.0\%$ error-target reconstruction yielded $15{,}396$ physical elements ($+10.4\%$ match to the reported $\approx 13{,}941$ elements). Authoritative production solve \texttt{1400395.mmaster02} completed $7{,}028$ increments with zero cutbacks and zero warnings (\texttt{Exit status 0}, walltime $07\,\mathrm{h}\,06\,\mathrm{min}$, CPUT $06\,\mathrm{h}\,52\,\mathrm{min}$), yielding $K_0 = 137.8437\,\mathrm{kN/mm}$ ($-0.08\%$), $F_{\text{peak}} = 0.748197\,\mathrm{kN}$ ($-1.29\%$), $u_{\text{peak}} = 0.005775\,\mathrm{mm}$ ($-1.45\%$), and a $98.31\%$ post-peak load drop, passing all five predeclared Task-5 scientific acceptance gates (Task 5: \textbf{QUANTITATIVE REPRODUCTION PASSED}). An accompanying controlled $5.0\%$ sensitivity solve (\texttt{1400396.mmaster02}, $4{,}194$ elements, walltime $01\,\mathrm{h}\,55\,\mathrm{min}$, CPUT $01\,\mathrm{h}\,50\,\mathrm{min}$) yielded $F_{\text{peak}} = 0.764964\,\mathrm{kN}$ ($+0.92\%$) and $u_{\text{peak}} = 0.007060\,\mathrm{mm}$ ($+20.48\%$), demonstrating that coarser crack-corridor refinement diffuses the phase damage gradient and delays peak localization.

Fourth, visualization integration (Task 6) was investigated. In the absence of accessible authentic IMFD ABAQUSER software, an in-solver companion facsimile UMAT visualization bridge was executed under production Job \texttt{1400408.mmaster02} (walltime $07\,\mathrm{h}\,23\,\mathrm{min}$, CPUT $07\,\mathrm{h}\,08\,\mathrm{min}$). The bridge maps UEL phase and history variables through Fortran \texttt{COMMON /CB\_STATE\_TRANS/} and \texttt{UEXTERNALDB} lifecycle logic into passive \texttt{CPE4}/\texttt{CPE3} UMAT elements, writing $\text{STATEV}(15) = d$ and $\text{STATEV}(16) = \mathcal{H}$ directly to the standard Abaqus output database. The calculation demonstrated exact $0.000000\%$ mechanical parity with baseline \texttt{1400395.mmaster02} across all $7{,}028$ increments. Evaluated phase-field outputs revealed a bounded numerical overshoot ($\max(d) = 1.00518084$), which was confirmed to originate in the solved UEL phase field rather than the visualization transfer (\textbf{\texttt{FORMULATION\_LEVEL\_NOT\_VISUALIZATION\_TRANSFER}}, classified as \textbf{\texttt{POSSIBLE\_FORMULATION\_DISCRETIZATION\_CAUSE}}). The companion bridge is \textbf{FULLY VERIFIED}, while formal integration of authentic ABAQUSER remains designated as \textbf{\texttt{TASK6\_BLOCKED\_EXTERNAL\_ABAQUSER\_DEPENDENCY}} pending software provision by the institute; Task 6 is therefore not scientifically complete.

Finally, nonmatching state transfer and multi-resolution restart mechanics were investigated, demonstrating runtime state ingestion while identifying residual mechanical-continuity challenges upon phase release across changing mesh topologies. The current shared-state UEL architecture is qualified strictly for serial execution due to mutable common storage. With Task 5 accepted and Task 6 verified at the companion-bridge level, the investigation advances to Task 7 benchmark and sensitivity studies.
